\documentclass[a4paper, serif, 10pt]{moderncv}

%% moderncv
% color and style
\moderncvcolor{black}
\moderncvstyle{banking}

%% page layout/margins
\usepackage[top=2.5cm, bottom=2.5cm, left=1.5cm, right=1.5cm]{geometry}

\nopagenumbers{}

%% Babel config for French language
% ref:
% - https://latex3.github.io/babel/guides/locale-french.html
\usepackage[french]{babel}

%% fontspec
\usepackage{fontspec}

\IfFontExistsTF{Linux Libertine}{
  \setmainfont[Ligatures={TeX, Common}, Numbers=OldStyle]{Linux Libertine}
}{}
\IfFontExistsTF{Linux Libertine O}{
  \setmainfont[Ligatures={TeX, Common}, Numbers=OldStyle]{Linux Libertine O}
}{}

%% CTeX
% CJK support, used to show CJK characters in the resume
%
% - heading=false: disable ctex heading
% - punct=kaiming: use kaiming punctuations styles for CJK
% - scheme=plain: use plain scheme, do not override \normalsize font size
% - space=auto: space settings for CJK characters
%
% CJK fonts are resolved per locale language, e.g. Noto Serif/Sans CJK JP
% for Japanese resumes, with fallback to the next region if not installed
%
% ref:
% - http://ctan.mirrorcatalogs.com/language/chinese/ctex/ctex.pdf
\usepackage[UTF8, heading=false, punct=kaiming, scheme=plain, space=auto]{ctex}

\IfFontExistsTF{Noto Serif CJK SC}{
  \setCJKmainfont{Noto Serif CJK SC}
}{}
\IfFontExistsTF{Noto Sans CJK SC}{
  \setCJKsansfont{Noto Sans CJK SC}
}{}

%% line spacing
\usepackage{setspace}
\setstretch{1.125}

%% URL styling - use normal text instead of monospace
\urlstyle{same}

%% Patch \httplink and \httpslink to handle full URLs
% The original moderncv macros always prepend http:// or https:// to URLs.
% This causes issues when the URL already has a protocol (e.g., https://example.com)
% resulting in malformed URLs like https://https://example.com.
% This patch checks if the URL already contains "://" and uses it directly if so.
% Note: We use \str_set:Nx (x-expansion) to expand macros like \@homepage first.
\ExplSyntaxOn
\str_new:N \g_yamlresume_url_str

\RenewDocumentCommand{\httpslink}{O{}m}{
  \str_set:Nx \g_yamlresume_url_str {#2}
  \str_if_in:NnTF \g_yamlresume_url_str {://}
    { \url{#2} }
    { \url{https://#2} }
}

\RenewDocumentCommand{\httplink}{O{}m}{
  \str_set:Nx \g_yamlresume_url_str {#2}
  \str_if_in:NnTF \g_yamlresume_url_str {://}
    { \url{#2} }
    { \url{http://#2} }
}
\ExplSyntaxOff

\name{Camille Laurent}{}
\title{Ingénieur en cybersécurité – Sécurité cloud et réponse à incident}
\phone[mobile]{+33 6 45 78 12 96}
\email{camille.laurent@protonmail.com}
\homepage{https://www.camille-laurent.fr}

\address{14 rue de la Paix, Paris, Île-de-France, France, 75002}{}{}

\extrainfo{{\small \faLinkedin}\ \href{https://www.linkedin.com/in/camille-laurent-cyber}{@camille-laurent-cyber} {} {} {} • {} {} {} 
{\small \faGithub}\ \href{https://github.com/camille-laurent}{@camille-laurent}}

\begin{document}

\maketitle

\section{Informations de base}

\cvline{}{Ingénieur en cybersécurité avec plus de huit ans d'expérience dans la protection des infrastructures critiques, la sécurité cloud et la réponse à incident. Reconnu pour concevoir des architectures Zero Trust, automatiser les opérations de sécurité et accompagner les équipes dans une démarche de conformité durable.
%
\begin{itemize}
\item Expertise en sécurité cloud AWS/Azure, architecture Zero Trust et DevSecOps
\item Pilotage d'opérations SOC, réponse à incident et analyse forensique
\item Maîtrise des référentiels ISO 27001, NIST CSF, RGPD et EBIOS RM
\item Encadrement technique, mentorat et communication auprès des parties prenantes
\end{itemize}}
\section{Formation}

\cventry{sept. 2014–juin 2016}
        {Master, Cybersécurité et systèmes d'information, Score : Mention Bien}
        {École Supérieure d'Ingénierie Numérique de Paris}
        {\url{https://www.esin-paris.fr}}
        {}
        {\begin{itemize}
\item Formation solide en sécurité des systèmes, cryptographie et gestion des risques
\item Projets pratiques sur la détection d'intrusion et la sécurisation d'applications web
\item Mémoire de fin d'études consacré à la sécurité des environnements cloud multi-tenant
\end{itemize}
\textbf{Cours} : Sécurité des réseaux et protocoles, Cryptographie appliquée, Analyse forensique, Gestion des risques et conformité, Sécurité des systèmes d'exploitation, Développement sécurisé, Virtualisation et sécurité cloud, Droit du numérique et protection des données}
\section{Expérience professionnelle}

\cventry{mars 2022–Present}
        {Ingénieur en cybersécurité senior – Sécurité cloud}
        {FinSecure Paris}
        {\url{https://www.finsecure-paris.fr}}
        {}
        {\begin{itemize}
\item Pilotage de la stratégie de sécurité cloud pour un groupe bancaire de 4 000 collaborateurs
\item Conception d'architectures Zero Trust sur AWS et Azure, avec automatisation des contrôles par IaC
\item Réduction de 35 \% des incidents critiques grâce au déploiement d'une plateforme SOAR et à la chasse aux menaces
\item Encadrement d'une équipe de quatre ingénieurs et animation du comité de sécurité applicative
\item Mise en conformité RGPD, ISO 27001 et NIST CSF, en lien avec les auditeurs externes
\end{itemize}
\textbf{Mots-clés} : Sécurité cloud, Zero Trust, SOAR, RGPD, ISO 27001, Mentorat}

\cventry{janv. 2019–févr. 2022}
        {Analyste SOC et réponse à incident}
        {CloudNova}
        {\url{https://www.cloudnova.fr}}
        {}
        {\begin{itemize}
\item Surveillance 24/7 des infrastructures cloud et qualification des alertes SIEM pour 60 clients
\item Investigation et réponse aux incidents de sécurité, de la détection à la remédiation
\item Mise en place de playbooks SOAR et enrichissement des alertes via le renseignement sur les menaces
\item Réalisation d'analyses forensiques et rédaction de rapports post-incident
\item Participation à des tests d'intrusion internes et à des revues de configuration
\end{itemize}
\textbf{Mots-clés} : SOC, SIEM, SOAR, Forensique, Réponse à incident, Tests d'intrusion}

\cventry{sept. 2016–déc. 2018}
        {Administrateur sécurité réseaux}
        {Groupe SantéAxe}
        {\url{https://www.santeaxe.fr}}
        {}
        {\begin{itemize}
\item Administration des équipements de sécurité réseau : pare-feu, proxy, IDS/IPS et VPN
\item Durcissement des systèmes et application des politiques de contrôle d'accès
\item Supervision des vulnérabilités et coordination des campagnes de correctifs
\item Sensibilisation des utilisateurs à l'hygiène numérique et aux risques de phishing
\end{itemize}
\textbf{Mots-clés} : Pare-feu, IDS/IPS, VPN, Durcissement, Vulnérabilités, Sensibilisation}
\section{Langues}

\cvline{Français}{Niveau natif ou bilingue}
\cvline{Anglais}{Niveau professionnel complet \hfill \textbf{Mots-clés} : TOEIC 940}
\cvline{Espagnol}{Niveau de travail limité \hfill \textbf{Mots-clés} : Niveau B1}
\section{Compétences}

\cvline{Sécurité des réseaux}{Avancé \hfill \textbf{Mots-clés} : Pare-feu, IDS/IPS, VPN, Segmentation, 802.1X}
\cvline{Sécurité cloud}{Expert \hfill \textbf{Mots-clés} : AWS, Azure, GCP, Zero Trust, CASB, CSPM}
\cvline{Opérations de sécurité}{Avancé \hfill \textbf{Mots-clés} : SIEM, SOAR, EDR, XDR, MITRE ATT\&CK, Renseignement sur les menaces}
\cvline{Gouvernance et conformité}{Avancé \hfill \textbf{Mots-clés} : ISO 27001, NIST CSF, RGPD, EBIOS RM, SOC 2}
\cvline{Sécurité applicative}{Avancé \hfill \textbf{Mots-clés} : OWASP Top 10, SAST, DAST, DevSecOps, IaC}
\cvline{Programmation et automatisation}{Avancé \hfill \textbf{Mots-clés} : Python, Bash, PowerShell, Terraform, Ansible, Docker, Kubernetes}
\section{Récompenses}

\cventry{nov. 2023}
        {École Supérieure d'Ingénierie Numérique de Paris}
        {Prix d'excellence en cybersécurité}
        {}
        {}
        {\begin{itemize}
\item Récompense décernée pour un projet de détection automatisée des menaces dans les environnements cloud
\item Distinction attribuée parmi 120 candidats pour l'impact opérationnel et la qualité de la recherche
\end{itemize}}
\section{Certificats}

\cventry{mars 2024}
        {ISC2}
        {CISSP – Professionnel certifié en sécurité des systèmes d'information}
        {\url{https://www.isc2.org/Certifications/CISSP}}
        {}
        {}

\cventry{juin 2023}
        {Amazon Web Services}
        {Certification AWS Security – Spécialité sécurité cloud}
        {\url{https://aws.amazon.com/certification/certified-security-specialty/}}
        {}
        {}

\cventry{sept. 2021}
        {CompTIA}
        {CompTIA Security+ – Certification en cybersécurité}
        {\url{https://www.comptia.org/certifications/security}}
        {}
        {}

\cventry{avr. 2022}
        {PECB}
        {ISO/IEC 27001 Lead Implementer – Mise en œuvre d'un SMSI}
        {\url{https://pecb.com/en/education-and-certification-for-individuals/iso-iec-27001/iso-iec-27001-lead-implementer}}
        {}
        {}
\section{Publications}

\cventry{févr. 2024}
        {Détection des intrusions par apprentissage automatique en environnement cloud}
        {Revue Française de Cybersécurité}
        {\url{https://www.rfcyber.fr/articles/detection-intrusions-apprentissage-automatique}}
        {}
        {\begin{itemize}
\item Étude comparative de modèles supervisés pour la détection d'anomalies réseau dans des environnements multi-cloud
\item Proposition d'un pipeline de réduction du bruit et d'enrichissement des alertes pour les équipes SOC
\item Retours d'expérience sur le déploiement en production et mesure du taux de faux positifs
\end{itemize}}
\section{Projets}

\cventry{janv. 2023–déc. 2023}
        {Plateforme open source d'automatisation de la qualification des alertes et de la réponse aux incidents.}
        {Plateforme SOAR de réponse à incident}
        {\url{https://github.com/camille-laurent/soar-incident-response}}
        {}
        {\begin{itemize}
\item Développement de playbooks Python pour l'enrichissement et la remédiation automatique
\item Intégration avec SIEM, EDR et solutions de ticketing via API REST
\item Réduction du temps moyen de réponse de 40 \% sur un périmètre de 200 serveurs
\end{itemize}
\textbf{Mots-clés} : SOAR, Python, Automatisation, API REST, SIEM}
\section{Centres d'intérêt}

\cvline{Cybersécurité offensive}{CTF, Chasse aux vulnérabilités, Rétro-ingénierie}
\cvline{Sport}{Escalade, Course à pied}
\cvline{Musique}{Guitare, MAO}
\section{Bénévolat}

\cventry{sept. 2019–juin 2024}
        {Mentor et formateur en cybersécurité}
        {Association CyberSensibilisation France}
        {\url{https://www.cybersensibilisation.fr}}
        {}
        {\begin{itemize}
\item Animation d'ateliers de sensibilisation auprès de collégiens, lycéens et demandeurs d'emploi
\item Conception de supports pédagogiques sur le phishing, les mots de passe et la protection des données
\item Accompagnement de jeunes professionnels dans leur préparation aux métiers de la cybersécurité
\end{itemize}}
    
\end{document}